% Exact definitions for the compact methodology. See methodology-notes.md.
% No new experimental protocol or result is asserted in this appendix.
\section{Intervention definitions}
\label{app:interventions}

\subsection{Sites and gate semantics}

Pythia uses parallel attention and FFN branches with separate LayerNorms.
For batch size $B$, sequence length $T$, residual width $d$, FFN width $d_f$,
and $H$ heads, let $d_h=d/H$. Table~\ref{tab:sites} identifies the actual
operands recorded after any selected gate. The same site mapping is applied
in every layer. A0 retains the stock block; A1-H replaces GELU at $h$ with
ReLU. A4 and A7 use the mappings in Section~\ref{sec:interventions}, with one
common threshold per condition. Pressure targets are specified separately;
the ladder's L1/OL1 suffixes apply pressure to all gated sites in that row.

\begin{table}[ht]
\centering\small
\begin{tabularx}{\linewidth}{@{}lllX@{}}
\toprule
Site & Shape & Consuming operation & Placement \\
\midrule
$a$ & $B\!\times\!T\!\times\!d$ & QKV projection & After attention LayerNorm \\
$m$ & $B\!\times\!T\!\times\!d$ & FFN up projection & After FFN LayerNorm \\
$h$ & $B\!\times\!T\!\times\!d_f$ & FFN down projection & Nonlinearity output \\
$q,k$ & $B\!\times\!H\!\times\!T\!\times\!d_h$ & Attention scores & After partial RoPE \\
$v$ & $B\!\times\!H\!\times\!T\!\times\!d_h$ & Probability--value & After the QKV split \\
$z$ & $B\!\times\!T\!\times\!d$ & Attention output & Concatenated context, before output projection \\
\bottomrule
\end{tabularx}
\caption{Intervention sites and the matrix products that consume them directly.}
\label{tab:sites}
\end{table}

Both threshold gates retain equality and leave retained values unchanged.
The comparison mask is detached: the derivative with respect to the input
is one for retained entries and zero for rejected entries. Consequently,
$G_{+,0}$ and ReLU have identical forward values but differ at zero, where
their derivatives are one and zero, respectively. $G_{\pm,0}$ is exactly the
identity in both values and gradients. The gates use neither a straight-through
estimator nor a top-$k$ selection rule. At an ungated $h$, the activation is
GELU; at other ungated sites it is the identity applied to that port's input.

\subsection{Activation pressure and OL1}
\label{app:pressure}

For the $J$ targeted site--layer tensors $A_j$, containing $n_j$ entries each,
the pressure objective is
$\mathcal{L}_1=J^{-1}\sum_{j=1}^{J}n_j^{-1}\sum_{i=1}^{n_j}|A_{j,i}|$.
Absolute values and means are evaluated in float32 on post-gate activations.
Each tensor receives equal weight regardless of its size; adding pressure
sites changes both the objective and the weight of existing targets.
Task and pressure losses are averaged over the same accumulation microbatches.
For naive L1, the accumulated gradient of
$\mathcal{L}_{\mathrm{task}}+\lambda\mathcal{L}_1$ is globally norm-clipped
before AdamW~\citep{loshchilov2019adamw}.

For OL1, let $g_t$ and $g_1$ be the accumulated task and unweighted pressure
gradients. We first clip $g_t$ by its global norm and perform AdamW using only
that clipped task gradient. Let $\widehat m_t$ and $\widehat v_t$ be its newly
updated, bias-corrected moments. On parameters having both gradients and
initialized optimizer state, define the adaptive task and pressure vectors
$u=\widehat m_t/(\sqrt{\widehat v_t}+\epsilon_A)$ and
$w=g_1/(\sqrt{\widehat v_t}+\epsilon_A)$. The pressure gradient is unclipped
and has no first-moment buffer. Inner products and norms below pool all these
eligible parameters. With $c=\langle u,w\rangle$, $q=\|u\|_2^2$, and numerical
stabilizer $\epsilon>0$, the conflict-conditioned projection is
\begin{equation}
\widetilde w=
\begin{cases}
w-\dfrac{c}{q+\epsilon}u, & c<0\ \text{and}\ q>\epsilon,\\[3pt]
w, & \text{otherwise}.
\end{cases}
\label{eq:ol1-projection}
\end{equation}
For pressure weight $\lambda\geq0$ and required budget $b>0$, set
$r=\lambda\|\widetilde w\|_2/(\|u\|_2+\epsilon)$ and
$s=\min\{1,b/(r+\epsilon)\}$ when $r>0$, or $s=1$ otherwise. After AdamW,
each eligible parameter group $p$, with learning rate $\alpha_p$, receives
$\theta_p\leftarrow\theta_p-\alpha_p\lambda s\widetilde w_p$.
The experiments' clipping thresholds, $b$, and $\lambda$ are condition
parameters; the reference clipping norm is one.

This procedure is related to gradient-conflict projection~\citep{yu2020pcgrad}
and the use of smoothed primary directions~\citep{hsieh2024bloop}, but defines
a particular correction in AdamW-preconditioned coordinates. The stabilized
projection is only approximately orthogonal under conflict. Its geometry and
budget precede multiplication by group learning rates and exclude decoupled
weight decay. They therefore do not guarantee task-loss descent or preservation.

\section{Sparsity accounting and architectural reach}
\label{app:accounting}

\subsection{Exact counters}

Local exact sparsity at site $s$ is
$\mathcal{Z}_s=\sum_{\ell,b,i}\mathbf{1}\{A^{(\ell,b)}_{s,i}=0\}/
\sum_{\ell,b}|A_s^{(\ell,b)}|$, where $\ell$ indexes layers, $b$ evaluation
batches, and $i$ tensor entries. Near-zero mass replaces equality by
$|A_{s,i}|\leq\varepsilon$ for a stated tolerance; it does not contribute to
exact-zero counts unless the value is zero. All aggregate fractions pool
integer numerators and denominators across batches and layers before division.

The atomic unit of model-wide accounting is one scalar multiplication in a
matrix product. If $X\in\mathbb{R}^{n\times p}$ feeds a linear projection with
$q$ outputs, the count is $N=npq$, with $N_0=q\sum_{i,j}\mathbf{1}\{X_{ij}=0\}$.
Zero weights receive no credit. In attention, both operands are activations,
so a product is credited once if either is zero. For query position $t$,
key position $u\leq t$, and within-head coordinate $j$,
\begin{align}
N_{0,QK}&=\sum_{b,h,t}\sum_{u\leq t}\sum_j
\mathbf{1}\{Q_{bhtj}=0\ \lor\ K_{bhuj}=0\},\label{eq:qk-count}\\
N_{0,PV}&=\sum_{b,h,t}\sum_{u\leq t}\sum_j
\mathbf{1}\{P_{bhtu}=0\ \lor\ V_{bhuj}=0\}.\label{eq:pv-count}
\end{align}
Here $P$ is the causal softmax probability tensor. Future masked positions
are excluded from both counts, while zeros on valid positions, including
floating-point underflow, are included. Query and key counts use actual
post-RoPE operands: rotation can repopulate an individual zero coordinate.
Likewise, the output projection uses actual post-gate context $z$, rather than
inferring its zeros from the value mask. Earlier value positions participate
in more causal pairs, so local value sparsity alone does not determine $PV$
sparsity.

\subsection{Denominator and ceiling derivation}
\label{app:ceiling}

For one uncached, unpadded sequence, each block contributes $3Td^2$ products
from QKV projection, $Td^2$ from the attention-output projection, $Td d_f$ from
each FFN projection, and $dT(T+1)/2$ from each of $QK$ and $PV$.
Thus, for $L$ blocks and vocabulary size $V_{\!\mathrm{vocab}}$,
\begin{equation}
\begin{split}
C_{\mathrm{block}}&=T(4d^2+2dd_f)+dT(T+1),\\
C_{\mathrm{model}}&=LC_{\mathrm{block}}+TdV_{\!\mathrm{vocab}}.
\end{split}
\label{eq:denominator}
\end{equation}
Over $B$ sequences, $N_{\mathrm{model}}=BC_{\mathrm{model}}$. The numerator
$N_0$ sums the six block families above, whereas the dense head contributes
only to the denominator. Embedding lookup, normalization, nonlinearities,
softmax, RoPE, bias and residual additions, and data movement are outside this
accounting. For a block-only view, $\Sblock=N_0/(BLC_{\mathrm{block}})$.
Both sparsities are fractions; percentages multiply them by 100.

Let $\mathcal{O}(\mathcal{G})$ be the union of operations reachable when the
selected sites are all zero: $a$ reaches QKV; $m$ reaches FFN up; $h$ reaches
FFN down; either $q$ or $k$ reaches $QK$; $v$ reaches $PV$ and attention output;
and $z$ reaches attention output. Each operation is counted once.
The $v$ rule includes the exact closure $V=0\Rightarrow PV=0$. No further
propagation is assumed: QKV biases can survive $a=0$, and zero scores do not
imply zero softmax probabilities. For the same site set in every layer,
$N_{\mathrm{reach}}=BL\sum_{o\in\mathcal{O}(\mathcal{G})}C_o$, where $C_o$
is that operation's per-sequence count. Division by $N_{\mathrm{model}}$ gives
$\Smax$.

With Pythia's $d_f=4d$, the per-block reachable counts for A0, A1-H, A4, and A7
are $0$, $4Td^2$, $12Td^2$, and $C_{\mathrm{block}}$, respectively. The ladder
evaluates these formulas at $T=2{,}048$, $V_{\!\mathrm{vocab}}=50{,}304$, and
the pinned $(L,d)$ configurations $(6,128)$, $(6,512)$, and $(24,1{,}024)$.
No checkpoint enters this calculation. In particular, A0 has zero selected-site
reach but may have positive observed sparsity from natural zeros. The ceiling
is consequently not an upper bound on every observed $\Smodel$, nor does its
gap to an observation measure sparsity available without quality loss.

\subsection{Evaluation scope}
\label{app:evaluation}

Quality and logical counts are paired in the same eager-attention evaluation
of a checkpoint, with its declared gates and evaluation precision. The complete
MiniPile validation split contains 500 documents. Concatenation and packing
produce 338 complete blocks of 2,048 tokens (692,224 input tokens); the final
1,444-token tail is excluded. Each block contributes 2,047 next-token targets
to cross-entropy, while product accounting includes all 2,048 input positions.
Loss is averaged over the equally sized blocks, and sparsity is count-pooled.
This workload uses no key--value cache. Condition-specific training budgets,
optimizer settings, seeds, and checkpoint selection accompany their results;
they are not implied by the ladder's list of candidate conditions.
